%
% 24-koerper.tex -- Körperkoeffizienten
%
% (c) 2026 Prof Dr Andreas Müller
%
\subsection{Polynome mit Koeffizienten in einem beliebigen Körper}
Der euklidische Algorithmus, wie er in
Abschnitt~\ref{buch:ratint:euklid:subsection:polynome}
für Polynome mit Koeffizienten in $\mathbb{Q}$ oder $\mathbb{C}$
vorstellt wurde, ist in keiner Weise auf solche Zahlkörper beschränkt.
Damit die Polynomdivision und damit die Bestimmung des Rests
in jedem Schritt des euklidischen Algorithmus durchgeführt werden kann,
ist nur nötig, dass die Koeffizienten aus einem Körper stammen.
Die Diskussion der Eigenschaften euklidischer Ringe in
Abschnitt~\ref{buch:ratint:euklid:subsection:ringe} hat dies
gezeigt.

Ist $K$ ein Körper, dann ist auch die Menge $F=K(x)$ der rationalen Funktionen
in der Variablen $x$ ein Körper.
$F$ ist ein euklidischer Ring.
Für Polynome mit Koeffizienten in $F$ gilt daher ebenfalls der
euklidische Algorithmus.

Eine Funktion wie
\[
f(z)
=
\frac{2z\log z + 9}{\frac1z-\log z}
\]
kann als rationale Funktion in der Variablen $\vartheta=\log z$
mit Koeffizienten im Körper $\mathbb{Q}(z)$ der rationalen Funktionen
in der Variablen $z$ betrachtet werden.
Zähler und Nenner sind Polynome
\[
2z\vartheta + 9
\qquad\text{und}\qquad
\frac1z -\vartheta
\]
in $\vartheta$ mit rationalen Funktionen als Koeffizienten.
Die Schreibweise
\[
f(z)
=
\frac{2z\vartheta - 9}{\frac1z-\vartheta}
\]
mit $\vartheta$ macht dies deutlich.
Wir werden dies im Kapitel~\ref{chapter:intalg} wiederholt ausnutzen,
um einen Integrationsalgorithmus zu konstruieren.
Wir werden dazu den Grundkörper der rationalen Funktionen in $z$ 
durch weitere Funktionen wie $\vartheta$ erweitern.
Die Erweiterung verlangt nur, dass wir mit rationalen Funktionen
umgehen können müssen.
Zwar ergeben sich zusätzliche Erschwernisse dadurch, dass die Koeffizienten
nicht konstant sind, aber die prinzipielle Vorgehensweise von
Abschnitt~\ref{buch:ratint:bernoulli:subsection:prinzip}
bleibt gültig.



